\documentclass[a4paper,11pt]{article}
\usepackage[T1]{fontenc}
\usepackage[utf8]{inputenc}
\usepackage{lmodern,microtype}
\usepackage[a4paper,margin=1in,headheight=15pt,headsep=14pt]{geometry}
\usepackage{amsmath,amssymb,graphicx,booktabs}
\usepackage[table,svgnames]{xcolor}
\usepackage[most]{tcolorbox}
\tcbuselibrary{skins,breakable}
\usepackage{pifont}
\IfFileExists{bbding.sty}{\usepackage{bbding}\newcommand{\glyphHand}{\Pointinghand}\newcommand{\glyphPencil}{\Pencil}}{\newcommand{\glyphHand}{\ding{43}}\newcommand{\glyphPencil}{\ding{46}}}
\usepackage{enumitem}
\setlist{leftmargin=1.5em,topsep=4pt,itemsep=3pt}
\newlist{steps}{enumerate}{1}
\setlist[steps]{label=\textbf{Step \arabic*.},leftmargin=4.4em,labelwidth=3.8em,labelsep=0.6em,labelindent=0pt,align=left,topsep=5pt,itemsep=6pt}
\usepackage{parskip,fancyhdr,titlesec}
\usepackage[hidelinks]{hyperref}
\usepackage{xurl}
\hypersetup{breaklinks=true}
\definecolor{navy}{HTML}{21355E}\definecolor{navyrule}{HTML}{21355E}\definecolor{inktext}{HTML}{1A202C}\definecolor{mutedtext}{HTML}{6B7280}
\definecolor{intuFrame}{HTML}{2C5AA0}\definecolor{intuFill}{HTML}{F4F7FC}
\definecolor{workFrame}{HTML}{2E7D52}\definecolor{workFill}{HTML}{F1F8F3}
\definecolor{keyFrame}{HTML}{6A4C93}\definecolor{keyFill}{HTML}{F4F0F9}
\color{inktext}
\hypersetup{linkcolor=navy,urlcolor=navy,citecolor=navy}
\newcommand{\CourseShort}{Machine Learning}
\newcommand{\SessionNo}{7}
\newcommand{\LectureTitle}{Decision Tree Classifier \& Information Theory}
\pagestyle{fancy}\fancyhf{}
\fancyhead[L]{\footnotesize\color{mutedtext}\textsf{\CourseShort}}
\fancyhead[R]{\footnotesize\color{mutedtext}\textsf{Session \SessionNo\ \textperiodcentered\ \LectureTitle}}
\fancyfoot[L]{\footnotesize\color{mutedtext}\textsf{WILP, BITS Pilani}}
\fancyfoot[C]{\footnotesize\color{mutedtext}\thepage}
\renewcommand{\headrulewidth}{0.4pt}\renewcommand{\footrulewidth}{0pt}
\titleformat{\section}{\sffamily\Large\bfseries\color{navy}}{\thesection}{0.8em}{}[{\vspace{2pt}\color{navy!70}\titlerule[0.8pt]}]
\titleformat{\subsection}{\sffamily\large\bfseries\color{navy!92}}{\thesubsection}{0.7em}{}
\titlespacing*{\section}{0pt}{16pt}{8pt}\titlespacing*{\subsection}{0pt}{12pt}{4pt}
\newif\ifmlkfirstsec \mlkfirstsectrue
\newcommand{\sectionbreak}{\ifmlkfirstsec\mlkfirstsecfalse\else\clearpage\fi}
\newcommand{\secsub}[1]{\noindent{\sffamily\bfseries\itshape\color{navy!85}\large #1}\par\medskip}
\tcbset{housecallout/.style 2 args={enhanced, breakable, colback=#2, colframe=#1, boxrule=0.6pt, arc=2.4pt, outer arc=2.4pt, left=10pt, right=10pt, top=9pt, bottom=9pt, before skip=11pt, after skip=11pt, fontupper=\normalsize, coltitle=white, fonttitle=\bfseries\sffamily\small, attach boxed title to top left={xshift=6mm,yshift=-3mm}, boxed title style={colback=#1, colframe=#1, rounded corners, arc=3pt, boxrule=0pt, left=6pt, right=6pt, top=2pt, bottom=2pt}}}
\newtcolorbox{intuition}{housecallout={intuFrame}{intuFill}, title={\raisebox{-0.05ex}{\glyphHand}\,\ The intuition}}
\newtcolorbox{worked}[1]{housecallout={workFrame}{workFill}, title={\raisebox{-0.05ex}{\glyphPencil}\,\ Worked example: #1}}
\newtcolorbox{keytake}{housecallout={keyFrame}{keyFill}, title={\raisebox{-0.05ex}{\ding{51}}\,\ Key takeaway}}
\newcommand{\bitswordmark}{{\sffamily\bfseries\color{white}\fontsize{12.5}{15}\selectfont WILP, BITS Pilani}}
\newcommand{\titlebanner}[4]{\begin{tcolorbox}[enhanced, breakable=false, colback=navy, colframe=navy, boxrule=0pt, arc=3pt, outer arc=3pt, left=16pt, right=16pt, top=16pt, bottom=16pt, before skip=2pt, after skip=14pt, width=\linewidth]\noindent\begin{minipage}[c]{0.72\linewidth}\bitswordmark\end{minipage}\hfill\par\vspace{9pt}{\sffamily\footnotesize\bfseries\color{white!85}\MakeUppercase{#1}}\par\vspace{6pt}{\sffamily\bfseries\fontsize{26}{30}\selectfont\color{white} #2\par}\vspace{5pt}{\sffamily\large\color{white!88} #3\par}\vspace{9pt}{\color{white!55}\rule{\linewidth}{0.4pt}}\par\vspace{6pt}{\sffamily\footnotesize\color{white!82} #4\par}\end{tcolorbox}}
\newcommand{\housefig}[2]{\begin{figure}[htbp]\centering\includegraphics[width=\linewidth]{#1}\caption{#2}\end{figure}}
\newcommand{\vocab}[1]{\textbf{\color{navy}#1}}
\begin{document}
\thispagestyle{fancy}
\titlebanner{Machine Learning \textperiodcentered\ Companion Reader}{Decision Tree Classifier \& Information Theory}{Hunt's Algorithm, Entropy, Information Gain, and ID3 Construction}{Session 7 \textperiodcentered\ Read alongside the lecture slides \textperiodcentered\ Every slide example worked out in full}
\smallskip\noindent{\small\color{mutedtext}Prepared for the Work Integrated Learning Programmes (WILP), BITS Pilani.}\smallskip
\section{First Taste: Twenty Questions for Data}
\secsub{A decision tree splits complex decisions into a sequence of simple yes/no questions.}
A \vocab{decision tree} converts tabular data into a flowchart structure. Each internal node tests an attribute, each branch represents an outcome, and each leaf node holds a class prediction.
\begin{intuition}
A decision tree breaks down a complex classification problem into a hierarchy of simple feature tests.
\end{intuition}
\housefig{figures/fig_decision_tree_structure.png}{Structure of a decision tree showing decision nodes, feature branches, and pure leaf node predictions.}
\section{Decision Tree Representation \& Expressiveness}
\secsub{Decision trees express disjunctions of conjunctions on attribute constraints.}
Any decision tree can be rewritten as a set of logical rules.
\subsection{Hunt's Algorithm}
Given dataset $D_t$ at node $t$:
\begin{enumerate}
  \item If all records in $D_t$ belong to class $y_k$, node $t$ becomes a leaf labeled $y_k$.
  \item If $D_t$ contains multiple classes, select an attribute test condition to partition $D_t$ into child subsets and recurse.
\end{enumerate}
\section{Information Theory \& Entropy}
\secsub{Entropy measures the amount of impurity, disorder, or unexpected surprise in a dataset.}
\begin{equation}
H(S) = -\sum_{i=1}^{c} p_i \log_2(p_i).
\end{equation}
\housefig{figures/fig_entropy_curve.png}{Binary entropy curve H(S) as a function of positive class proportion p+.}
\begin{worked}{Computing Entropy for Sample Class Distributions}
Calculate binary entropy for sample subsets:
\begin{steps}
  \item \textbf{Pure dataset $[14+, 0-]$:} $H(S) = 0$ bits.
  \item \textbf{Balanced dataset $[7+, 7-]$:} $H(S) = 1.0$ bit.
  \item \textbf{Slightly imbalanced dataset $[9+, 5-]$:} $H(S) \approx 0.940$ bits.
\end{steps}
\end{worked}
\section{Information Gain \& The ID3 Algorithm}
\secsub{Information Gain measures the expected reduction in entropy from splitting on an attribute.}
\begin{equation}
\text{Gain}(S, A) = H(S) - \sum_{v \in \text{Values}(A)} \frac{|S_v|}{|S|} H(S_v).
\end{equation}
\housefig{figures/fig_info_gain_bars.png}{Information Gain values for candidate attributes at the root node.}
\begin{worked}{14-Sample PlayTennis Root Split Selection}
In the PlayTennis dataset ($S = [9+, 5-]$, $H(S) = 0.940$ bits):
\begin{steps}
  \item \textbf{Evaluate Outlook:} $\text{Gain}(S, \text{Outlook}) = 0.940 - 0.694 = 0.246$ bits.
  \item \textbf{Compare Root Candidates:} Outlook ($0.246$), Humidity ($0.151$), Wind ($0.048$), Temp ($0.029$). Outlook is chosen as root.
\end{steps}
\end{worked}
\end{document}
